\section{Modelagem Matemática}

Esta seção especifica os AFDs que reconhecem os padrões de placa
definidos na especificação da linguagem.

\subsection{Notação e alfabeto}

\begin{center}
\begin{tabular}{cl}
\toprule
Símbolo & Definição \\
\midrule
$\Lt$ & as 26 letras maiúsculas, $\{A, B, \dots, Z\}$ \\
$\Nm$ & os algarismos, $\{0, 1, \dots, 9\}$ \\
$\Hf$ & o conjunto $\{\texttt{-}\}$, contendo o hífen \\
$\varepsilon$ & palavra vazia (ausência do hífen na forma antiga) \\
\bottomrule
\end{tabular}
\end{center}

O alfabeto comum aos autômatos é $\Sigma = \Lt \cup \Nm \cup \Hf$.
O hífen pertence ao alfabeto, embora seja rejeitado pelo autômato
Mercosul. As letras são consideradas apenas em caixa alta; a
normalização de minúsculas, se desejada, ocorre antes da simulação.

Nas tabelas, $\Lt$, $\Nm$ e $\Hf$ representam qualquer símbolo da
classe correspondente. Toda transição não indicada leva ao estado
de rejeição $q_R$, e $\delta(q_R, a) = q_R$ para todo $a \in \Sigma$.

% ------------------------------------------------------------------
\subsection{AFD 01: padrão anterior ao Mercosul (PNU)}

O padrão anterior tem três letras seguidas de quatro algarismos. Na
chapa física, o hífen pode aparecer depois das três letras; em
registros de dados, pode estar ausente. A linguagem é, portanto,
$\Lt^3\,\Hf?\,\Nm^4$.

\subsubsection{Definição formal}

$M_1 = (Q_1, \Sigma, \delta_1, q_0, F_1)$, em que
\begin{align*}
  Q_1 &= \{q_0, q_1, q_2, q_3, q_4, q_5, q_6, q_7, q_8, q_R\},\\
  F_1 &= \{q_8\}.
\end{align*}

Os estados $q_1$, $q_2$ e $q_3$ registram as três letras; $q_4$
registra a leitura do hífen; $q_5$, $q_6$, $q_7$ e $q_8$ registram os
quatro algarismos; $q_R$ é o estado de rejeição.

\subsubsection{Função de transição}

\begin{table}[h]
\centering
\begin{tabular}{cccc}
\toprule
$\delta_1$ & $\Lt$ & $\Nm$ & $\Hf$ \\
\midrule
$\to q_0$ & $q_1$ & $q_R$ & $q_R$ \\
$q_1$     & $q_2$ & $q_R$ & $q_R$ \\
$q_2$     & $q_3$ & $q_R$ & $q_R$ \\
$q_3$     & $q_R$ & $q_5$ & $q_4$ \\
$q_4$     & $q_R$ & $q_5$ & $q_R$ \\
$q_5$     & $q_R$ & $q_6$ & $q_R$ \\
$q_6$     & $q_R$ & $q_7$ & $q_R$ \\
$q_7$     & $q_R$ & $q_8$ & $q_R$ \\
$* q_8$   & $q_R$ & $q_R$ & $q_R$ \\
$q_R$     & $q_R$ & $q_R$ & $q_R$ \\
\bottomrule
\end{tabular}
\caption{Função de transição do AFD 01.}
\label{tab:delta1}
\end{table}

A partir de $q_3$, a leitura de um algarismo leva diretamente a
$q_5$, o que permite a forma sem hífen. A leitura de $\Hf$ leva a
$q_4$, que exige um algarismo em seguida. Assim, o hífen é opcional
e, quando presente, aparece exatamente após a terceira letra.

\subsubsection{Linguagem reconhecida}

\[
  \lang{PNU} = \{\, xhy \mid x \in \Lt^3,\; h \in \{\varepsilon, \texttt{-}\},\; y \in \Nm^4 \,\}
\]

São aceitas \texttt{ABC-1234} e \texttt{ABC1234}. São rejeitadas
\texttt{AB-12345}, \texttt{ABC--1234}, \texttt{ABC-12A4} e
\texttt{ABC1-234}.

% ------------------------------------------------------------------
\subsection{AFD 02: padrão Mercosul}

O padrão Mercosul tem a sequência LLLNLNN: três letras, um
algarismo, uma letra e dois algarismos. Neste reconhecedor geral, a
quinta posição aceita qualquer letra de A a Z.

\subsubsection{Definição formal}

$M_2 = (Q_2, \Sigma, \delta_2, q_0, F_2)$, em que
\begin{align*}
  Q_2 &= \{q_0, q_1, q_2, q_3, q_4, q_5, q_6, q_7, q_R\},\\
  F_2 &= \{q_7\}.
\end{align*}

Os estados $q_1$, $q_2$ e $q_3$ registram as três letras iniciais;
$q_4$ registra o primeiro algarismo; $q_5$ registra a letra da quinta
posição; $q_6$ e $q_7$ registram os dois algarismos finais.

\subsubsection{Função de transição}

\begin{table}[h]
\centering
\begin{tabular}{cccc}
\toprule
$\delta_2$ & $\Lt$ & $\Nm$ & $\Hf$ \\
\midrule
$\to q_0$ & $q_1$ & $q_R$ & $q_R$ \\
$q_1$     & $q_2$ & $q_R$ & $q_R$ \\
$q_2$     & $q_3$ & $q_R$ & $q_R$ \\
$q_3$     & $q_R$ & $q_4$ & $q_R$ \\
$q_4$     & $q_5$ & $q_R$ & $q_R$ \\
$q_5$     & $q_R$ & $q_6$ & $q_R$ \\
$q_6$     & $q_R$ & $q_7$ & $q_R$ \\
$* q_7$   & $q_R$ & $q_R$ & $q_R$ \\
$q_R$     & $q_R$ & $q_R$ & $q_R$ \\
\bottomrule
\end{tabular}
\caption{Função de transição do AFD 02.}
\label{tab:delta2}
\end{table}

\subsubsection{Linguagem reconhecida}

\[
  \lang{Mercosul} = \Lt^3 \cdot \Nm \cdot \Lt \cdot \Nm^2
\]

São aceitas \texttt{BRA2E19}, \texttt{QTP9Z01} e \texttt{ABC1D23}.
São rejeitadas \texttt{AB12C34}, \texttt{ABC12D3}, \texttt{ABC1D2A}
e \texttt{ABC-1D23}.

% ------------------------------------------------------------------
\subsection{AFD 03: Mercosul convertido (extensão opcional)}

Como extensão, a quinta posição pode ser restrita às letras de A
até J, conjunto $\LAJ = \{A, B, C, D, E, F, G, H, I, J\}$. A
estrutura continua sendo LLLNLNN.

\subsubsection{Definição formal}

$M_3 = (Q_3, \Sigma, \delta_3, q_0, F_3)$, com $Q_3 = Q_2$ e
$F_3 = \{q_7\}$. A função $\delta_3$ coincide com $\delta_2$, exceto
nas transições que partem de $q_4$:
\[
  \delta_3(q_4, a) =
  \begin{cases}
    q_5 & \text{se } a \in \LAJ,\\
    q_R & \text{se } a \in \Lt \setminus \LAJ.
  \end{cases}
\]

\subsubsection{Linguagem reconhecida}

\[
  \lang{Conv} = \Lt^3 \cdot \Nm \cdot \LAJ \cdot \Nm^2
\]

Vale $\lang{Conv} \subseteq \lang{Mercosul}$. O AFD verifica apenas a
restrição formal da quinta posição; ele não determina a origem ou o
histórico real de uma placa.

% ------------------------------------------------------------------
\subsection{Reconhecimento combinado}

O programa principal executa os dois reconhecedores obrigatórios de
forma independente: primeiro o AFD 01 e, se ele rejeitar a entrada,
o AFD 02. A placa é aceita se pelo menos um deles a aceitar:
\[
  \lang{Validador} = \lang{PNU} \cup \lang{Mercosul}.
\]
A versão restrita (AFD 03) pode ser usada como verificação adicional,
sem substituir o reconhecedor Mercosul geral.

\begin{table}[h]
\centering
\begin{tabular}{lccl}
\toprule
Entrada & AFD 01 & AFD 02 & Resultado \\
\midrule
\texttt{ABC-1234} & aceita  & rejeita & válida (anterior) \\
\texttt{ABC1234}  & aceita  & rejeita & válida (anterior) \\
\texttt{BRA2E19}  & rejeita & aceita  & válida (Mercosul) \\
\texttt{QTP9Z01}  & rejeita & aceita  & válida (Mercosul) \\
\texttt{AB12C34}  & rejeita & rejeita & inválida \\
\texttt{ABC12D3}  & rejeita & rejeita & inválida \\
\texttt{ABC1D2A}  & rejeita & rejeita & inválida \\
\bottomrule
\end{tabular}
\caption{Classificação de entradas pelos dois reconhecedores.}
\label{tab:classificacao}
\begin{definicao}
Um autômato finito determinístico é uma tupla
$M = (Q, \Sigma, \delta, q_0, F)$, em que $Q$ é o conjunto de
estados, $\Sigma$ o alfabeto, $\delta : Q \times \Sigma \to Q$ a
função de transição, $q_0$ o estado inicial e $F \subseteq Q$ o
conjunto de estados finais \cite{hopcroft2002}.
\end{definicao}

\emph{Em palavras:} $Q$ são as etapas da leitura, $\Sigma$ são os
caracteres permitidos, $\delta$ diz para qual etapa ir ao ler cada
caractere, $q_0$ é onde a leitura começa e $F$ são as etapas em que
a placa é considerada válida.

\emph{Em palavras:} o AFD 01 exige três letras, aceita um hífen
opcional e termina em quatro algarismos. Qualquer outro
caractere, em qualquer ponto, leva ao estado de erro $q_R$.
\end{table}